\section{Introduction}

The nocturnal stable boundary layer (SBL) remains one of the most persistent unresolved regimes in boundary-layer meteorology. Under strong stratification, local closure-based scaling frameworks become progressively less reliable as a full description of observed transition behavior. Across field and model studies, the same forcing classes can produce sharply different outcomes: sustained turbulent coupling in some intervals and abrupt decoupling in others \citep{mcnider1993, mcnider1995, vandewiel2017, vignon2017}. This mismatch motivates a shift in emphasis from single-point threshold interpretation toward state-space structure.

Historically, the SBL paradox is visible in repeated reversed-S organization of Richardson-based diagnostics, with branch-like behavior and hysteretic switching \citep{mcnider1993, mcnider1995}. A local-metric reading treats this as scatter or non-stationary variability around closure curves. A geometric reading treats the same pattern as projection overlap: multiple physically distinct column states mapping to similar local coordinates. In that view, the central issue is not whether Richardson diagnostics are useful, but where they sit in the coordinate hierarchy. They are retained here as local indicators embedded within a broader manifold representation of atmospheric-column evolution.

This study develops that representation using reduced coordinates $(\eta_1,\eta_2,\eta_3)$ derived from p-FEM profile reconstruction and metric-consistent projection. We apply the framework across CASES-99, FLOSS, BLLAST, and GABLS3 to test whether folded topology persists under strong contrasts in terrain, forcing context, and measurement geometry. The goal is not to replace classical diagnostics, but to place them in a chart-equivalent setting where local thresholds can be interpreted as projections of global state-surface motion.

The Results section is organized as an explicit evidence sequence. First, we establish shared folded geometry across campaigns (Figure~\ref{fig:manifold}). Second, we characterize brittle versus rubbery transition archetypes as campaign-dependent fold-edge behavior within the same structural class (Figure~\ref{fig:archetypes}). Third, we show projection shadowing in time series, where local $\Rig$ can saturate while $\eta_3$ remains dynamically informative (Figure~\ref{fig:shadow}). Fourth, we quantify fold-proximal precursor structure through lag and slowing diagnostics (Figure~\ref{fig:slowing}). Fifth, we evaluate transferability through curvature-coordinate collapse against fold-distance scaling (Figure~\ref{fig:collapse}).

Throughout, causal claims are constrained to what is directly measured in each campaign or simulation context. This guardrail is methodological, not stylistic: it separates direct observations from interpretation-level hypotheses and prevents over-extension across field and LES domains while preserving a unified geometric framework.
